\documentclass[11pt]{article}
\usepackage{fontspec}
\usepackage{amsmath,amssymb}
\usepackage[margin=1in]{geometry}
\usepackage{url}
\begin{document}
\sloppy
\section*{GROK H1427 (2026-10-06 14:27 PT): independent re-simulation and 1e6 seed tally}
\textbf{P versus NP: OPEN.}

\section*{0. What landed since H1411}
\texttt{main} tip is \texttt{057c50b0} (H1411). Two peer commits landed just before it, \texttt{cac9f744} and \texttt{7c32aa4f} (TWOSATDEC28, $q=33$ exact-layer draw), under \path{results/raw/monotone_slice/} and \path{research/explorations/monotone-slice/}; neither is under \path{research/agents/grok/} or \path{research/main/}, and H1411 missed mentioning them. Nothing new under the watched paths.

\section*{1. Literature}
None. The next claim concerns this loop's own control circuit.

\section*{2. Independent re-simulation of the control circuit (COMPUTATIONAL)}
H1411's check of the eight-gate control circuit used the same truth-table code as the solver encoding. This pass re-simulated it with a standalone script (\path{/workspace/scratch_h1427_1006/indep.py}, sha256 \texttt{c7f35781}) that shares no code with \path{enum8peb.py}. Under the convention that input $i$ is bit $i$ of the row index and gate table bit $q$ is the output on $(a,b)=(q\bmod 2,\lfloor q/2\rfloor)$, the circuit $(6,0,4,14)$, $(7,3,6,6)$, $(8,2,6,8)$, $(9,0,5,6)$, $(10,1,8,14)$, $(11,9,10,8)$, $(12,7,11,14)$, $(13,8,12,6)$ outputs exactly \texttt{a40ff54e580f5d0a}. Swapping operand order also matches, because tables 6, 8 and 14 are symmetric. Reversing the input order gives \texttt{a05fa00f72be703c}, which is not $f$, so the match depends on the convention and is not vacuous. This confirms the circuit computes $f$. It does not re-check the two-pebble length 11, which still rests on the shared BFS.

\section*{3. Seeds at $10^6$ conflicts (COMPUTATIONAL, partial)}
At 14:29\,PT, 20 of the 74 non-control seeds are done (10 per worker), all UNKNOWN, each in 84.7 to 100.9\,s; 0 SAT, 0 UNSAT. The other 54 are still running (about 91\,s each per worker, so roughly 40 more minutes). The control was SAT in 55\,s at the same budget, but UNKNOWN still does not mean UNSAT. The machine has 8 cores at load about 9, so no new job was started this pass.

\section*{4. Pending run read: tight $L=10$ (COMPUTATIONAL, partial)}
\path{l10.out} at 14:28\,PT: cubes 0 to 29 hit the 900\,s cap; 0 SAT, 0 UNSAT, 30 timeouts out of 75 kept cubes. Timeouts are not UNSAT. Whether $L(f,2)$ is 10 or 11 is OPEN.

\section*{5. Standing status}
No claim is falsified and no charged resource improves this pass. $10\le L(f,2)\le 11$ stands, the lower bound CONDITIONAL on the tightness lemma (CANDIDATE), the upper bound COMPUTATIONAL. The largest unconditional confirmed price stays one. The time-state-representation surface stays CANDIDATE. The SAT-to-charged-model compiler and the product-span obstruction are both OPEN. The D(4,8) semantic-cut bound of 48 stays FALSIFIED/RETIRED.

\section*{Next smallest step}
Write a standalone two-register pebbling checker that replays the H1411 program step by step against the definition, so the length 11 no longer rests on the shared BFS. Keep reading \path{run1m_*.out} and \path{l10.out}; if all 74 seeds end UNKNOWN, stop per-seed budgets as H1411 planned.

\bigskip
\noindent\textbf{Status sentence.} P versus NP remains OPEN: the SAT-to-charged-model compiler and the product-span obstruction are open load-bearing arrows. This pass independently re-simulates the H1411 control circuit and confirms it computes \texttt{a40ff54e580f5d0a}, finds 20 of 74 other first-gate seeds UNKNOWN at $10^6$ conflicts (0 SAT, 0 UNSAT), and finds tight $L=10$ cubes 0 to 29 still timed out (0 SAT, 0 UNSAT); $10\le L(f,2)\le 11$ is unchanged, with the lower bound CONDITIONAL on the tightness lemma. No claim is falsified, and this is not a proof of anything about P versus NP.
\end{document}
